\section{A simple example script}
\label{sec:example}

Figure~\ref{fig:example-script} gives a simple (and somewhat contrived)
example script, intended to introduce the language.  The script expects to
find integers in an initial segment of column~|A|.  It calculates their
factorials, and writes them in column~|B|.  It also prints the total of those
factorials below.  

We describe the script below.  The script is executed in order.

%%%%%

\begin{figure}
\begin{scala}[numbers = left]
// The input and output columns.
val In = #A; val Out = #B £\label{line:InOut}£

// The factorial of £n£.
def fact(n: Int): Int = { assert(n >= 0); if(n <= 1) 1 else n*fact(n-1) }£\label{line:factorial}£

// The sum of £xs£.
def sum(xs: List[Int]): Int = if(xs == []) 0 else head(xs) + sum(tail(xs))£\label{line:sum}£

// Row for the first empty cell in column £c£ from row £r£.
def firstEmptyInCol(c: Column, r: Row): Row = £\label{line:firstEmptyInCol}£
  Cell(c,r) match{ case Empty => r; case n: Int => firstEmptyInCol(c, r+1) }£\label{line:cell-match}£

// Row for the first empty cell in £In£
val firstEmpty = firstEmptyInCol(In, #0)£\label{line:firstEmpty}£

// Write factorials of entries in £In£ into £Out£.
for(r <- #0 until firstEmpty) Cell(Out,r) = fact(Cell(In,r))£\label{line:write-facts}£

// Write the total.
Cell(In, firstEmpty+1) = "Total:"£\label{line:total}£
Cell(Out, firstEmpty+1) = sum([Cell(Out,r): Int | r <- #0 until firstEmpty])£\label{line:write-sum}£
\end{scala}
\caption{A simple example script.}
\label{fig:example-script}
\end{figure}

%%%%%

The characters ``|\\|'' Turn the rest of the line into a comment. 

Line~\ref{line:InOut} defines two values |In| and |Out|, representing the
input and output columns.  It is a good idea to give names to relevant rows
and columns, in the same way that it is a good idea to give names to relevant
constants in standard programming languages.  A column value can be defined by
preceding the corresponding letter with ``|#|'', so, for example, ``|#A|''
represents column~|A|.  Likewise a row value can be defined by preceding the
row number with ``|#|'', for example ``|#0|''.  We call these forms \emph{row}
and \emph{column literals}.

Line~\ref{line:factorial} defines the factorial function.  Definitions of
functions are introduced using the keyword |def|.  Each parameter of a
function is given a type.  The return type of a function should be given when
the function is recursive---which is the case for all three functions in this
script---or when a reference is made to the function before its definition
(e.g.~when two functions are defined by mutual recursion).  The |assert|
statement is an assertion, capturing the natural precondition of the factorial
function.  The rest of the definition of |fact| is a straightforward
recursion, whose meaning should be obvious.

The |sum| function at line~\ref{line:sum} calculates the sum of a list of
|Int|s.  If |A| is a type, then |List[A]| is the type of lists that contain
values of type~|A|.  The value |[]| is the empty list.  More generally, a list
can be defined by writing its elements between square brackets, separated by
commas, for example ``|[2, 4, 6, 8]|''.  The function |head| returns the first
element of a nonempty list; and the function |tail| returns the list
containing all except the first element of a nonempty list.  Again, the
definition is a straightforward recursion.  

The expression |Cell(c,r)| represents the value in the spreadsheet cell in
column~|c| and row~|r|.  In the case where |c| and |r| are column and row
literals, this can be written more concisely.  For example, |#B3| is shorthand
for |Cell(#B,#3)|.  The script may make assumptions about the types of values
that it reads from cells; if a cell does not contain the expected type, an
error is generated.

The function |firstEmptyInCol(c, r)|, at line~\ref{line:firstEmptyInCol},
finds the first empty cell in column~|c|, from row |r| onwards, and returns
the relevant row.  |Column| is the type of columns, and |Row| is the type of
rows.  The function requires the cell read to be either empty or an integer.
This is checked using pattern matching on the value, using syntax similar to
that in Scala.  Each case in a pattern matching expression is of the form
|case|~$pat$~|=>|~$e$, where $pat$ is a pattern, and $e$ is an expression
giving the result when that pattern is matched.  The pattern |Empty| matches
an empty cell.  Hence, if |Cell(c,r)| is empty, |firstEmptyInCol| returns~|r|.
The pattern |n: Int| matches an integer, and binds the name~|n| to that
integer in the subsequent expression.  Hence, if |Cell(c,r)| contains an
integer, the |firstEmptyInCol| function recurses with
\SCALA{firstEmptyInCol(c, r+1)}, where |r+1| represents the row after~|r|;
thus, the function simply continues searching from the next row.  
%% We could have written the pattern as ``|_: Int|'' in this case, where
%% ``|_|'' represents a wildcard pattern, because the subsequent expression
%% does not use the value of the integer.
If the value read from the cell does not match any of the patterns---in this
case, if it is not empty or an integer---the pattern matching fails, and the
spreadsheet application gives an error message.

Line~\ref{line:firstEmpty} sets |firstEmpty| to be the first row in column
|In| that contains an empty cell. 

A cell can be written using a command of the form ``|Cell(c,r) = |$e$'',
where~$e$ defines the value written.  For example, line~\ref{line:total}
writes the string ``Total:'' into the cell below the first empty cell in
column~|In|.

Line~\ref{line:write-facts} writes into column~|Out| the factorials of values
in corresponding cells in column~|In|.  It does this using a |for| loop.  A
|for| loop of the form |for(x <- xs)|~$block$ executes $block$ for each value
of~|x| from the list~|xs|.  Here $block$ is a block containing declarations
and commands.  The expression |#0 until firstEmpty| represents the list of
rows from~|0|, inclusive, until |firstEmpty|, exclusive, i.e.~the rows that
were previously found to contain integers in column~|In|.  For each such
row~|r|, line~\ref{line:write-facts} writes the corresponding factorial into
column~|Out|.

Finally, line~\ref{line:write-sum} calculates the sum of the factorials, and
writes it into the appropriate cell.  The expression in square brackets on the
right is a list comprehension, using Haskell-style notation.  It forms the
list of values in row~|r| of column~|Out|, for each value of~|r| in \SCALA{#0
  until firstEmpty}, i.e.~the factorials written at
line~\ref{line:write-facts}.  The subexpression |Cell(Out, r): Int| tells
the spreadsheet application that the cell should contain an integer; the type
checker makes use of this information.
